\begin{table}[htbp]
  \centering
  \caption{Profile of formal workers by firm-size group, 2025}
  \label{tab:formal-worker-profile-2025}
  \small
  \begin{tabular}{lrrr}
    \toprule
    Characteristic & Formal solo & Formal 2--100 & Formal 101+ \\
    \midrule
    Share of formal workers & 6.4\% & 38.7\% & 54.9\% \\
    Mean real hourly income & 15,990 & 12,696 & 16,944 \\
    Mean age & 44.9 & 38.3 & 37.9 \\
    Women & 43.3\% & 41.9\% & 45.7\% \\
    Higher education & 50.9\% & 48.2\% & 63.4\% \\
    Own-account workers & 82.1\% & 7.7\% & 9.6\% \\
    Domestic workers & 16.2\% & 0.9\% & 0.0\% \\
    Professional and administrative services & 22.8\% & 14.6\% & 9.4\% \\
    Transport and storage & 15.6\% & 5.5\% & 5.6\% \\
    Commerce and repair & 14.7\% & 21.1\% & 10.1\% \\
    Health and social assistance & 5.2\% & 6.5\% & 10.6\% \\
    Households as employers & 16.2\% & 0.9\% & 0.0\% \\
    \bottomrule
  \end{tabular}
  \vspace{0.3em}
  \begin{minipage}{0.95\textwidth}
  \footnotesize
  Notes: Statistics use 2025 formal workers and GEIH expansion weights. The first row reports each column's share of all formal workers in the table. The 2--100 column pools all formal workers in firm-size categories from 2--3 through 51--100 workers. Higher education corresponds to workers classified as superior or university educated. Sector rows report the weighted share of each group located in the listed economic activity.
  \end{minipage}
\end{table}
